\cotefile{157-4}
\scanpage{164}
\mspage{\unsure{38|39}}

d'autre part l'hypothèse sur $W$ implique
que $W$ est stable par facteurs directs, \deleted{\unsure{donc}}
$C$ l'est aussi\,; $\unsure{A|R}$ est clos à gauche, donc $\inserted{TC{=}}C\cap W$
aussi, donc par prop. 4.1. cor. 2 on a
$TC=\widetilde{TC}$, $TC$ est clos à gauche. On gagne
donc, par la dernière assertion du cor. 3 prop. \unsure{4.3|5.3|3.3}. (p. 35).

En fait, indépendamment des conditions
particulières sur $\unsure{\mathcal M|M},W$, on vient de prouver
ceci :

\uline{Lemme} 4.4.1 (Cor. de prop. 4.3) Soit $\mathcal M$ une
catégorie, $W$ un localiseur dans $\mathcal M$ \unsure{multiplicativement|modérément}
saturé et stable par facteurs directs. Alors
les structures de Quillen disons $(C,F)$
associées à $W$ s'obtiennent ainsi : on
prend une partie \deleted{\unsure{du} $\unsure{C|G}$ \unsure{g}} $C_0$ de $\mathrm{Fl}\mathcal M$
telle que
\[
\deleted{\unsure{(4.4.2)}}\qquad TF\overset{\mathrm{déf}}{=}C_{0*}\subset W,
\]
on pose
\[
C=\widetilde C_0=TF^*,\quad TC=C\cap W,\quad F=(TC)_{\unsure{*}},
\]
et on \deleted{vérifie que} doit se débrouiller
pour que $(C,TF)$ et $(TC,F)$ satisfassent
à la condition de factorisation de
Quillen, i.e. soient des pré-couples de Quillen.

[En fait $C,TF\inserted{\text{ et }F}$ étant clos \deleted{\unsure{par}} \unsure{par} construction,
et $TC=C\cap W$ étant stable par facteurs directs
($C$ et $W$ l'étant), ce sont \unsure{$=$|là} des \uline{couples} de
Quillen.]

% archive sheet number: 163
\scanpage{165}
\mspage{\deleted{\unsure{2}}39}

Les hypothèses \inserted{\unsure{sur} $W$\unsure{(|C|,}} faites \unsure{tantôt} sur $\unsure{\mathcal M|M}$ \deleted{\unsure{assurent}} servent
à assurer que si je prends $C_0$ petit, \unsure{donc|alors} la
factorisation marche pour $(C,TF)$ automatiquement. Il reste le problème pour
$TC,F$.

Notons que si elle marche, alors $TC=\widetilde{TC}=F^*$,
donc $TC$ est clos à gauche, et a fortiori
satisfait les conditions de stabilité correspondantes
\[
\left\{\begin{aligned}
&\text{a) Stable par cochgt. de base, par composition, par}\\
&\qquad\text{facteurs directs (ces \inserted{deux} derniers sont assurés}\\
&\qquad\text{automatiquement, car $C,W$ sont stables),}\\
&\text{b) Stable par $\unsure{\underrightarrow{\lim}|\underline{\lim}}$ ordinales strictes. \deleted{Ce qui}}
\end{aligned}\right. \tag*{(4.4.3)}
\]
\deleted{\unsure{à cet égard on}} L'hypothèse \inserted{de stabilité} sur $W$ sert \unsure{(}entre autres) à garantir
\deleted{que} la stabilité b) pour $TC$ (en fait, il
suffirait pour cela \inserted{seul} que $W$ soit stable par
$\underrightarrow{\lim}$ ordinales strictes). Il reste la stabilité
de $TC=C\cap W$ par cochgt. de base, qui
est problématique, car $W$ n'est en général
\uline{pas} stable par cochgt. de base. C'est pour
assurer cette stabilité que je me vois
pratiquement contraint ici à imposer
\[
 C_0\subset\mathrm{Cof}_W,\tag*{(4.4.\deleted{\unsure{3}}\unsure{4|3})}
\]
je dis que sous les hypothèses de stabilité
faites sur $W$, cela implique
\[
 C\overset{\mathrm{déf}}{=}\widetilde C_0\subset\mathrm{Cof}_W.\tag*{(4.4.5)}
\]

% archive sheet number: 164
\scanpage{166}
\mspage{40}

En effet, on a

\uline{Proposition} 4.4.2. Soit $\unsure{\mathcal M|M}$ stable par \inserted{petites} $\unsure{\underrightarrow{\lim}|\underline{\lim}}$ inductives
filtrantes, et $W$ \inserted{$\subset\mathrm{Fl}\mathcal M$} un localiseur \inserted{(mult. saturé)}
\uline{stable par les dites limites}. Alors $\mathrm{Cof}_W$ satisfait les
conditions de stabilité (4.4.3) (``stabilité
de Quillen gauche''), et est \unsure{stable} aussi
par $\underrightarrow{\lim}$ filtrantes quelconques.

En effet, la condition a) est satisfaite
sans condition \unsure{suppl.} sur $W$. D'autre part,
il est immédiat que cette condition
\unsure{p.r|pour|par} aux $\underrightarrow{\lim}$ filtrantes implique que
$\mathrm{Cof}_W$ \uline{est stable par limites inductives
filtrantes} dans $\uline{\mathrm{Fl}}\,\mathcal M$.

\uline{Cor. 4.4.3} Sous les conditions de la prop.
4.4.\unsure{2|7}., $\mathrm{Cof}_W\cap W=W^{\mathrm{univ}}$ (\unsure{Ens.|ens.} des
$W$-équivalences co-universelles dans $\mathcal M$) est
stable par limites inductives filtrantes,
et \deleted{est donc est} donc ``Qstable à gauche''
au sens de (4.4.3.).

\uline{Corollaire}. 4.4.4. Sous les conditions de
% slanted
la prop. 4.4.2\deleted{\unsure{,}}*, \inserted{supposons de plus que les
sommes amalgamées finies existent dans
$\mathcal M$, et} \deleted{supposons de plus que} soit $C_0$ une \inserted{petite} partie de $\mathrm{Cof}_W$,
\inserted{dont les éléments soient accessibles, \deleted{et \unsure{coaccessibles}}}
\deleted{$\mathrm{Fl}\mathcal M$ \unsure{close|clos} à gauche. Alors $\unsure{C\cap W}$ satisfait}
\inserted{\unsure{et} considérons}
\[
 C=\widetilde C_0,\quad TC\overset{\mathrm{déf}}{=}C\cap W.
\]
Alors on a
\[
 \deleted{\unsure{F}}C\subset\mathrm{Cof}_W,\quad TC=C\cap W^{\mathrm{univ}}\subset W^{\mathrm{univ}},
\]

% archive sheet number: 165
\scanpage{167}
\mspage{41}

\deleted{En effet} et $C$ et $TC$ sont \inserted{Q-}stables à gauche.

En effet, par le th. 1 du \S\,1 (p. 10) et
son corollaire, $\widetilde C_0$ est la plus petite partie
de $\mathrm{Fl}\unsure{\mathcal M|M}$ contenant $C_0$ et Q-stable à gauche.
Comme $\mathrm{Cof}_W$ \deleted{$\unsure{\subset}$} est Q-stable à g. et contient
$C_0$, on a donc $C=\widetilde C_0\subset\mathrm{Cof}_W$. Donc
$TC\overset{\mathrm{déf}}{=}C\cap W\subset\mathrm{Cof}_W\cap W\overset{\mathrm{déf}}{=}W^{\mathrm{univ.}}$,
\deleted{D'autre part $TC$} donc $TC=C\cap W^{\mathrm{univ}}$. Comme
$W^{\mathrm{univ}}$ est \inserted{lui aussi} Q-stable à g. (cor. précédent), il
en est de $=$ de $C\cap W$, qed.

\uline{Corollaire} 4.4.5. Supposons que \deleted{tout} objet de
$\mathcal M$ soit accessible. Sous les conditions du
corollaire précédent, \deleted{$\unsure{TF=C_*|F=C_*|C_*}$} posant
\[
 TF=C_{0*},=C_*,\qquad \deleted{\unsure{TC=C_*}}\ F=(TC)_*,
\]
le couple $(C,TF)$ est un couple de
Quillen. Pour qu'il en soit de $=$
du couple $(TC,F)$, il suffit qu'il
existe une \uline{petite} partie $TC_0$ de $C\cap W$,
telle que $\widetilde{TC_0}=TC\;(\overset{\mathrm{déf}}{=}C\cap W)$.

La première assertion est, \inserted{on l'a vu tantôt,} le th. 1 du
\S\,1. \deleted{déjà} La deuxième résulte \unsure{de|du} ce $=$
th. On peut préciser dans ce sens :

\uline{Prop. 4.4.\unsure{6|5}} Soit $\mathcal M$ une catégorie stable par petites
$\underrightarrow{\mathrm{lim}}$ filtrantes, par sommes amalgamées, et dont
tous les objets sont accessibles, $W$ un localiseur

% archive sheet number: 166
\scanpage{168}
\mspage{42}

\unsure{grossièrement|modérément} saturé dans $\unsure{\mathcal M|M}$, stable par $\unsure{\underline{\lim}|\underrightarrow{\lim}}$ filtrantes,
$C_0$ une \uline{petite} partie de $\mathrm{Fl}\mathcal M$, \deleted{\unsure{contenue}} telle
que
\[
 C_0\subset\mathrm{Cof}_W,\qquad \underset{\substack{\Vert\\(TF)}}{C_{0*}}\subset W,
\]
\unsure{posons}
\[
 C=\widetilde C_0,\ TC=C\cap W,\ TF=C_{\unsure{*}},\ F=(TC)_*,
\]
et soit $TC_0$ \deleted{une petite} partie de $TC=C\cap W$.
On a alors \deleted{$\unsure{TC_0\subset TC}$, $\unsure{F_0\supset F}$}
\deleted{Dans ces conditions \unsure{pour tout} $C_0$, il y a équivalence :}
% margin note written sideways
\margin{$(4.4.6)\ \left\{\begin{aligned}
&C\subset\mathrm{Cof}_W,\ \deleted{\unsure{TC_0\subset C}}\ TC=C\cap W^{\mathrm{univ}}\quad\text{(pour mémoire)}\\
&\widetilde{TC_0}\subset TC,\quad F_0\supset F\quad\text{(\unsure{ce dernier|ce dernier est} \deleted{\unsure{lieu}} trivial)}
\end{aligned}\right.$}
\deleted{$\unsure{(4.4.6)}\quad\unsure{TC_0}\quad\widetilde{TC_0}\subset TC,\quad\unsure{(TC_0)_*}\ F_0\overset{\mathrm{déf}}{=}(TC_0)_*\supset F$,}
et les conditions suivantes sont équivalentes,
\[
\begin{aligned}
\mathrm{a})&\quad \widetilde{TC_0}=TC\ \deleted{\unsure{\subset}} \quad\text{(on a $\subset$)}\\
\mathrm{a}')&\quad \deleted{\unsure{TC_*=(TC_0)_*}}\ F_0=F\quad\text{(on a $\supset$)}\\
\mathrm{b})&\quad TF=F_0\cap W\quad\text{(on a $\subset$)}\\
\mathrm{c})&\quad \unsure{a')|c')}W=TF\circ\widetilde{TC_0}\quad\text{(on a $\supset$)}
\end{aligned}
\]
\margin{\uline{NB} On voit facilement que $\boxed{TF=F\cap W}$ mais \unsure{ce|ceci} ne suffit pas pour avoir $TF=F_0\cap W$ pour $TC_0$ \unsure{donné}.}
\margin{\uline{NB} On sait que $\boxed{W=TF\circ TC}$, ce qui n'est nullement suffisant pour impliquer c) !}
\deleted{[\unsure{d)}\inserted{\unsure{bis}}\ \unsure{$(W,C,TC)$ satisfait à factorisation}}
et impliquent que $(W,C,F)$ est un triple
\uline{de Quillen \unsure{clos}.}

\uline{Dém} L'inclusion $\widetilde{TC_0}\subset TC$ résulte de la Q-stabilité
% slanted
\inserted{à g. de $TC$ (\unsure{cor.} 4.4.4)
et du fait que $TC_0\subset TC$ est petite
(\unsure{à spécifier|à figurer} dans le th. !
tout comme pour $C_0\subset\mathrm{Cof}_W$ \ldots)}
\deleted{(4.4.6) est \unsure{prouvée}.}
\deleted{à g. de $TC$ (4.4.4) et du fait que $\unsure{C_0}\subset\mathrm{Cof}_W$}
L'équnce de a), b), c) n'est autre que
la prop. 4.3, appliquée aux $W$ et aux deux
couples de Quillen $(C,TF),(\widetilde{TC_0},F_0)$ (compte
tenu du th. 1 du \S\,1). \deleted{et} Qu'on a donc
une \unsure{structure|str} de Quillen close a été vu
dans le cor. 3 de prop. 4.3 (p. 35). \deleted{\unsure{Comme}}
Il reste à prouver l'équivalence de
a et de a'. On a $\mathrm a\Rightarrow\mathrm a'$, puisque
$(TC_0)_*=(\widetilde{TC_0})_*=F_0$. \deleted{Mais} Inversement,
a') implique $\widetilde{TC_0}=\widetilde{TC}$, donc le th. de
factorisation pour $(TC,F)$ en vertu du
th. de factorisation cor. 1 (p. 11) appliqué à
\deleted{couple} $TC_0\subset TC\subset\widetilde{TC_0}$ et à $F_0=F=(TC_0)_*$.

% archive sheet number: 167
\scanpage{169}
\mspage{43}

$\unsure{[}$Ainsi, le seul problème, une fois trouvé $C_0$
(petite) satisfaisant $C_0\subset\mathrm{Cof}_W,\ C_{0*}\subset W$,
c'est de trouver une \uuline{petite} $TC_0\subset TC$
qui ``engendre'' $TC$``$=C\cap W$'', i.e. qui satisfasse
aux conditions équivalentes a, a', b, c de la
proposition précédente. J'ai passé des jours
à essayer de prouver l'existence de $TC_0$,
en supposant $\unsure{\mathcal M|M}$ muni d'une filtration
\unsure{cardinal|cardinale}, \inserted{et} $W$ une partie accessible de $\mathcal M$,
sans y être arrivé. Bien sûr, \deleted{\unsure{tout}} ce qui
nous importe vraiment, c'est que $(TC,F)$ satisfasse
à la condition de factorisation. Mais je
ne vois aucun moyen de le prouver, si ce
n'est en ``engendrant'' \inserted{à g.} $TC$ \deleted{\unsure{à g.}} par une \uline{petite}
\unsure{sous-ensemble} $TC_0$.

5. \uline{Résultats techniques préliminaires, sur les parties
accessibles de $\mathrm{Ob}\,\mathcal M$, stables par $\underrightarrow{\mathrm{lim}}$ filtrantes}

Je crois pourtant que je peux arriver à
quelque chose d'utilisable, en travaillant systématiquement \unsure{avec|sur} des $C$ ($\subset\mathrm{Cof}_W$, Q-closes à
gauche, Q-engendrées \inserted{à g.} par une \uline{petite} partie
$C_0$ \deleted{\unsure{)}} de $C$) qui sont stables par \unsure{ttes|toutes} limites
inductives filtrantes.

Je suppose à nouveau \deleted{\unsure{\mathcal M}} sur $\mathcal M$ les
conditions suivantes :

% archive sheet number: 168
\scanpage{170}
\mspage{44}

a) $\unsure{\mathcal M|M}$ stable par petites limites inductives (et pas
seulement filtrantes)

b) $\mathcal M$ \deleted{``accessible'' (\unsure{objectivement}), i.e.} muni d'une
filtration cardinale $(\mathrm{Filt}^{\unsure{\pi|\varpi}})_{\pi\geq\pi_0}$

c) Les objets de $\mathcal M$ sont accessibles (NB. il suffit
de \unsure{le|la} supposer pour les objets de $\mathrm{Filt}^{\pi_0}$).

b) implique que $\mathcal M$ contient une petite
famille génératrice par épim. stricts (p. ex.
une sous-catégorie \inserted{$C$} petite de $\underline{\mathrm{Filt}}^{\pi_0}$ \inserted{$=\mathcal M^{\pi_0}$} telle
que l'inclusion $C\unsure{-|\to}\underline{\mathrm{Filt}}^{\pi_0}$ soit une
équivalence de catégories). Par SGA 4 I 8.12.7,
il en résulte que $\mathcal M$ est stable par petites
limites projectives. Il s'ensuit le

\uline{Lemme 5.1.} (Sous les conditions a), b) ) Pour
tout objet $X$ de $\mathcal M$, l'ensemble $\mathrm{Ss}(X)$
des sous-objets de $X$, muni de la relation
d'ordre naturelle ``d'inclusion'', \deleted{\unsure{\}|;}} admet des
$\mathrm{Inf}$ et des $\mathrm{Sup}$ quelconques

Des $\mathrm{Inf}$, car ceux-ci s'expriment comme
des limites projectives dans $\mathcal M$. Or l'existence
des $\mathrm{Inf}$ implique celle des $\mathrm{Sup}$,
puisque $\underset{J}{\mathrm{Sup}}=\underset{J'}{\mathrm{Inf}}$, où $J'$ est l'ens.
des majorants de $J$ dans $\mathrm{Ss}(X)$.

\uline{Cor. 5.2} : Pour \unsure{toute|\deleted{tout}} flèche $f:Z\longrightarrow X$
dans $\mathcal M$, il existe, parmi les sous-objets \inserted{$X'$} de $X$
par lesquels $f$ se factorise, un plus petit

% archive sheet number: 169
\scanpage{171}
\mspage{45}

(qu'on notera $\mathrm{Im}\,f$). Le morphisme
$Z\longrightarrow\mathrm{Im}\,f$ est un épimorphisme.

Cette dernière assertion \deleted{\unsure{revient}} provient de
celle-ci :

\uline{Lemme}. \unsure{5.?|5.3|5.7}. Soit \deleted{\unsure{$C$}} $\unsure{\mathcal M|M}$ une catégorie
\inserted{où} les noyaux de doubles flèches existent.
\inserted{Soit} $\unsure{g|f}:Z\longrightarrow Y$ une flèche de $\mathcal M$ \deleted{qui} \inserted{Alors $g$ est épi \unsure{ssi} $\unsure{\exists|]}$, $g$}
ne se factorise pas par un sous-objet \inserted{\uline{strict}} $Y_0$ de $Y$ différent de $Y$.
\deleted{Alors $g$ est un épimorphisme.}

\deleted{(Il suffit $=$ que $g$ ne se factorise pas
par un sous-objet \uline{strict} -- \unsure{notamment|autrement},
par les $\mathrm{Ker}(f_1,f_2)$ \ldots)}

On va faire une hypothèse supplémentaire
(pratiquement anodine) sur les
$\mathrm{Filt}^{\unsure{\pi|\varpi}}$\deleted{\unsure{,}} :

d) $\mathrm{Filt}^{\pi}$ est stable par limites inductives
indexées par des catégories d'indices $I$
telles que $\mathrm{card}^{\mathrm{fl}}(I)\leq\pi$\textsuperscript{*}. (Il suffit de
supposer la stabilité par objet initial, sommes
amalgamées simples, et $\underrightarrow{\mathrm{lim}}$ filtrantes indexées
par un ensemble ordonné $I$ de cardinal $\leq\pi$.)
% slanted
\margin{Cette condition sera nommée (\unsure{$\mathrm{d}_0$|d°}) par la suite,}

De plus, si $f:Z\longrightarrow X$ est une flèche
de $\mathcal M$ telle que \inserted{$Z$} \deleted{$\unsure{X|Z}$} $\in\mathrm{Filt}^{\pi}$, alors $\mathrm{Im}\,f\in\mathrm{Filt}^{\pi}$.

Cette dernière condition \inserted{sur $\mathrm{Im}\,f$} \deleted{résulte des}
\deleted{autres, quand les épimorphismes et les monomorphismes sont stricts,} se décompose en
deux

% archive sheet number: 170
\scanpage{172}
\mspage{46}

$\mathrm d_1$) Tout sous-objet d'un objet de $\mathrm{Filt}^{\unsure{\pi|\varpi}}$ est dans $\mathrm{Filt}^{\pi}$

$\mathrm d_2$) Si $Z$ est dans $\mathrm{Filt}^{\pi}$, alors \deleted{tout quotient}
$Z'$ de $Z$ tel que \deleted{l'image} $\mathrm{Im}\,p=Z'$ (où
$p:Z\to Z'$ est le morphisme canonique)
est dans $\mathrm{Filt}^{\pi}$.

% slanted
\margin{Il semblerait qu'on \uline{n'a} pas besoin de $\mathrm d_1$), seulement
de $\mathrm d_2$. Mais en fait, $\mathrm d_1$) doit être conséquence de
$\mathrm d_2$, et $\mathrm d_2$ est conséquence de la \unsure{première partie} de d
(stabilité par \uline{lim} finies) si \unsure{tout} \deleted{\unsure{il}} épimorphisme de $\unsure{\mathcal M|M}$
est strict, donc effectif (si les $\mathrm{Filt}^{\pi}$ sont stables
par \uline{lim} finies \unsure{aussi|-- si} \ldots).}

Bien sûr, \deleted{le plus souvent, si} pratiquement $\mathcal M$ sera
tel que $\mathrm{Filt}^{\pi}$ est stable par
passage aux sous-objets \deleted{quelconques} et aux
quotients d'un objet (pas seulement les
quotients particuliers dont il est question
dans $\mathrm d_2$). D'ailleurs, lorsque \deleted{tout}
bimorphisme (i.e. mono \uline{et} épi) de $\mathcal M$
est iso, alors $\mathrm d_2$ \deleted{implique} \inserted{équivaut} bien à la
stabilité de $\mathrm{Filt}^{\pi}$ par quotients d'objets.

\uline{Proposition} 5.3. Supposons les conditions
a) b) \unsure{d$_0$)} \deleted{\unsure{d)}} vérifiées pour $\mathcal M$. Soit
$\pi$ un cardinal $\geq\pi_0$, et $X$ un objet
de $\mathcal M$. \deleted{Alors} Soit $\mathrm{Ss}_{\pi}(X)$ l'ens.
des sous-objets de $X$ qui sont dans $\mathrm{Filt}^{\pi}$.
Cet ensemble, \unsure{ordonné|muni} par la rel. d'ordre
induite par $\mathrm{Ss}(X)$ est \deleted{filtrant}.
\unsure{gr.|grand} devant $\pi$ (et en particulier filtrant),
\deleted{\unsure{et lorsque} $X$ \unsure{est} \illegible{2} \unsure{isomorphisme}
\unsure{limite inductive} \illegible{1} $X$}
\deleted{\unsure{de plus}} si de plus, les limites
inductives filtrantes de monomorphismes
de $\mathcal M$ sont mono\textsuperscript{*} \inserted{(p. ex. les limites inductives
filtrantes exactes à g.)}, ou si les objets de $\mathrm{Filt}^{\pi}$
sont $\pi$-accessibles, on a
$\underset{\alpha\in\mathrm{Ss}_{\pi}(X)}{\underrightarrow{\mathrm{lim}}} Z_{\alpha}\overset{\sim}{\longrightarrow} X$.
% margin note written sideways
\margin{C'est une \uline{condition \deleted{\unsure{de}} d'exactitude} \unsure{de} $\mathcal M$,
que \unsure{je vais|je veux} noter comme \textcircled{\unsure{e}} \deleted{\unsure{Alors}} les \unsure{mono|monos} de $\mathcal M$ sont stables
par \uline{lim} filtrantes dans $\underline{\mathrm{Fl}}(\mathcal M)$.}
% slanted
\margin{Dire aussi que $\mathrm{Ss}_{\pi}(X)$ est $\pi$-\unsure{complet|adapté|cofiltrant}, et le système
induc\textsuperscript{\unsure{tif}} \unsure{(disp\textsuperscript{f} $\alpha$)} $(Z_\alpha)_\alpha$, \unsure{aussi}.}

% archive sheet number: 171
\scanpage{173}
\mspage{47}

\uline{Dém} Soit $(\deleted{\unsure{X}}Z_\alpha)_{\alpha\in A}$ une famille de sous-objets
de $X$ tels que $\deleted{\unsure{X}}Z_\alpha\in\mathrm{Filt}^{\unsure{\pi|\varpi}}$, et $\mathrm{card}\,A\leq\pi$.
Soit $Z=\coprod_{\alpha\in A}\underset{\unsure{,,}}{\deleted{\unsure{X}}Z_\alpha}$, on a $Z\in\mathrm{Filt}^{\pi}$ par d).
Le plus petit sous-objet de $X$ majorant les
$Z_\alpha$ n'est autre que $\mathrm{Im}(Z\to X)$, qui
est dans $\mathrm{Filt}^{\pi}$ par $\mathrm d_{\unsure{2|1}}$. Prouvons que
$\underset{\alpha\in\mathrm{Ss}_{\pi}(X)}{\underrightarrow{\mathrm{lim}}}\,\deleted{\unsure{X}}Z_{\alpha}\overset{\sim}{\longrightarrow}X$ est iso. \deleted{On sait qu'} \inserted{Quand on sait que c'est un mono}
\unsure{--|a} $X=\underset{I}{\underrightarrow{\mathrm{lim}}}\,X_i$, où $(X_i)_{i\in\unsure{J|I}}$ est
un système inductif filtrant, $I$ un
ordonné grand devant $\pi$, \inserted{les $X_i$ dans $\mathrm{Filt}^{\pi}$.} Soit, pour tout
$i\in I$, $X'_i$ l'image de $X_i$ dans $X$.
\inserted{On a donc $X'_i\in\mathrm{Filt}^{\pi}$ par $\mathrm d_2$).}
Utilisant le fait que les $X_i\to X'_i$ sont
épi, on voit de suite que $X=\underset{I}{\underrightarrow{\mathrm{lim}}}\,X'_i$.
Considérons l'application
\diagram{\begin{tikzcd}[ampersand replacement=\&,row sep=0.2em,column sep=2em]
 I \arrow[r] \& \mathrm{Ss}_{\pi}(X) \\
 i \arrow[r,mapsto] \& X'_i
\end{tikzcd}}
C'est une application croissante, \inserted{soit $J$ son image,
partie filtrante grande devant $\pi$ de $\mathrm{Ss}_{\pi}(X)$} \deleted{pour que l'} telle
\deleted{est cofinale. Comme $I,\unsure{I'|J'}$ sont des ens. ordonnés
filtrants, \unsure{tout revient} cela signifie que pour
tout $Z_\alpha$ dans $\mathrm{Ss}_{\pi}(X)$, existe un $i$ tel que
$Z_\alpha\subset X'_i$.} que $\underset{\alpha\in J}{\underrightarrow{\mathrm{lim}}}\,Z_\alpha\overset{\sim}{\longrightarrow}X$.

% archive sheet number: 172
\scanpage{174}
\mspage{48}

On a donc
\diagram{\begin{tikzcd}[ampersand replacement=\&,row sep=1.2em,column sep=2em]
 \underset{\alpha\in J}{\underrightarrow{\mathrm{lim}}}\,Z_\alpha \arrow[r] \arrow[d,phantom,"\unsure{\wr|{\cong}}" description] \&
 \underset{\alpha\in\mathrm{Ss}_{\unsure{\pi|\varpi}}X}{\underrightarrow{\mathrm{lim}}}\,Z_\alpha \arrow[r,"\varphi"] \& X,\\
 X \& \&
\end{tikzcd}}
le composé étant l'identité. Si on sait
que $\varphi$ est un mono, il en résulte
que c'est un \unsure{iso} (mono. ayant
une section). \deleted{\unsure{Or}} Or $\varphi$ \unsure{sera} un
mono si \unsure{tte|toute|\deleted{tout}} $\underrightarrow{\mathrm{lim}}$\unsure{;|:} filtrante de
mono. est mono, \inserted{et alors on gagne.}
\deleted{puisque $\varphi$ est la limite inductive des
$(Z_\alpha\xrightarrow{i_\alpha} X'_\alpha\inserted{=X})$ correspondant à l'hom\unsure{omorphisme}
évident du syst. inductif $(Z_\alpha)_{\alpha\in\mathrm{Ss}_{\pi}X}$ dans le
\unsure{const}.}

Si on suppose \deleted{\unsure{donc}} que les
éléments de $\mathrm{Filt}^{\pi}$ \unsure{ont|sont} $\pi$-\unsure{accessible|accessibles},
on gagne aussi, car on trouve

\uline{Corollaire 5.4.} Si \deleted{\unsure{les}} objets de $\pi$
sont $\pi$-accessibles (p. ex. si $\pi\geq\pi_1$ et $\pi=\pi^{\pi_0}$,
p. ex. $\pi\geq\pi_1$ et $\pi=2^c$, avec $c\geq\pi_0$, cf
(3.2.) p. 16), alors pour tte partie filtrante
$J$ de $\mathrm{Ss}(X)$ grande devant $\pi$, telle que
$\underset{\alpha\in J}{\unsure{\mathrm{lim}|\underrightarrow{\mathrm{lim}}}}\,Z_\alpha\overset{\sim}{\longrightarrow}X$,
$J$ est cofinale dans $\mathrm{Ss}_{\pi}X$,
(donc $\underset{\alpha\in J}{\underrightarrow{\mathrm{lim}}}\,Z_\alpha\overset{\sim}{\longrightarrow}
\underset{\alpha\in\mathrm{Ss}_{\pi}(X)}{\underrightarrow{\mathrm{lim}}}\,Z_\alpha$)\deleted{\unsure{)}}.
\deleted{En effet, \unsure{soit}} C'est immédiat.

% archive sheet number: 173
\scanpage{175}
\mspage{4\deleted{\unsure{8}}9}

\deleted{\uline{Corollaire} \unsure{5.5.} Dans la définition d'une filtration
cardinale (SGA 4 I 9.12. p. 144), quand on
fait l'hypothèse de stabilité supplémentaire
d) sur les $\mathrm{Filt}^{\unsure{\pi|\varpi}}$, \inserted{$+$} dans la condition c) de \unsure{loc. cit.}
on peut supposer \deleted{que} \inserted{dans} la représentation de
$X$ comme $\underset{i\in I}{\underrightarrow{\mathrm{lim}}}\,X_i$ ($X_i\in\mathrm{Filt}^{\pi}$), $I$ grand
devant $\pi$, on ait \unsure{simplemt}, $I=\mathrm{Ss}_{\pi}(X)$}

\uline{Scholie} 5.5. Quand on fait \inserted{l'hyp. d'exactitude assez}
anodine que les $\underrightarrow{\mathrm{lim}}$ filtrantes de mono sont
mono, \uline{ou} qu'on se borne \unsure{aux cardinaux|au cardinal}
$\pi$ \unsure{tels|telle} que les objets de $\mathrm{Filt}^{\pi}$ soient
$\pi$-accessibles, on peut \uline{canoniser} la
représentation d'un objet $X$ comme $\underset{i\in I}{\underrightarrow{\mathrm{lim}}}\,X_i$,
avec $X_i\in\mathrm{Filt}^{\pi}$ et \deleted{\unsure{les}} $I$ grand devant $\pi$,
en prenant $I=\mathrm{Ss}_{\pi}(X)$, et le système
inductif correspondant. \deleted{(si \uline{lim} filtrante de mono est mono)}
On a de plus l'agréable propriété supplémentaire que
voici : toute partie filtrante \inserted{$J$} de $\unsure{I|II}$,
de cardinal $\leq\pi$, a une borne
supérieure $\deleted{\unsure{i}}\,i_J$, et (\deleted{si} \deleted{\unsure{l'on}} \unsure{suppose|supposons} que
$\unsure{\mathcal M|M}$ satisfait la condition \unsure{e)} page \unsure{46|45}, \unsure{bas page)|puisque l'|toujours)} on
aura
\[
 X_{i_J}=\underset{i\in J}{\underrightarrow{\mathrm{lim}}}\,X_i,
\]
en d'autres termes, on a une représentation
\deleted{\unsure{can}} de $X$ par un syst. inductif \inserted{(parfaitement)}
\uline{$\pi$-adapté}.
\inserted{(cf $\overline{\underline{\mathrm{XIII}}}$ p. 85, 86)}

% archive sheet number: 174
\scanpage{176}
\mspage{50}

Soit maintenant
\[
 M\subset\mathrm{Ob}\,\unsure{\mathcal M|M}
\]
une partie de $\mathrm{Ob}\,\mathcal M$, stable par la relation d'isomorphie. Si $M$ est $\unsure{\pi|\varpi}$-accessible \deleted{\unsure{alors}} (suivant la
condition $\mathrm b_\pi$ de $\overline{\underline{\mathrm{XIII}}}$, p. \unsure{88}), alors \deleted{\unsure{pour}}
posant
\[
 M^\pi=M\cap\mathrm{Filt}^\pi,
\]
$M^\pi$ détermine $M$ par la condition suivante
\inserted{sur $X\in\unsure{\mathrm{ob}|\mathrm{Ob}}\,\mathcal M$}
\[
\begin{aligned}
 X\in \deleted{\unsure{\mathcal M}}M\ \Longleftrightarrow\ &\text{le sous ens. }\mathrm{Ss}^M_\pi(X)\text{ de }\mathrm{Ss}_\pi(X)\\
 &\text{formé des }X_\alpha\in\mathrm{Ss}_\pi(X),\text{ tels que }X_\alpha\in M\\
 &(\Longleftrightarrow X_\alpha\in M^\pi)\text{ est cofinal dans }\mathrm{Ss}^{\deleted{\unsure{M}}}_\pi(X).
\end{aligned}\tag*{(5.1.)}
\]
Notons que lorsque $X\in\mathrm{Filt}^\pi$, ce critère est
trivialement valable, puisque alors $\mathrm{Ss}_\pi(X)$ a
$X$ comme \deleted{objet final} \inserted{plus grand élément}, et les sous-ens.
cofinaux sont ceux qui contiennent
cet objet final. Donc (5.1) est intéressant
seulement pour les $X\in\mathrm{Ob}\,\mathcal M\setminus\mathrm{Filt}^\pi$.

Notons que pour \uline{\unsure{tte|toute}} partie $M^\pi$ de
$\mathrm{Filt}^\pi$, on peut définir une partie $M$
de $\mathrm{Ob}\,\mathcal M$ par la condition (5.1.) et on
aura trivialement
\[
 M^\pi=M\cap\mathrm{Filt}^\pi.\tag*{(5.2.)}
\]
Mais sans doute, bien que $M$ satisfasse
(5.1.), il n'est pas pour \deleted{\unsure{autant}} \inserted{nécessairement}
$\pi$-accessible. En effet, pour qu'il le
soit, il \unsure{faut|doit} qu'il soit stable par $\underrightarrow{\mathrm{lim}}$

% archive sheet number: 175
\scanpage{177}
\mspage{51}

grandes devant $\unsure{\pi|\varpi}$ d'éléments de $M^\pi$, ce qui
implique que \deleted{\unsure{tout}} facteur direct d'un
élément de $M^\pi$ est dans $M$, donc aussi
dans $M^\pi$ d'après la condition de stabilité
d) sur $\mathrm{Filt}^{\pi}$. C'est là une condition
restrictive sur la partie $M^\pi$ de $\mathrm{Filt}^{\pi}$.

Aussi, on devra avoir une condition
sur $\unsure{\mathcal M|M}$ plus forte que la validité de
(5.1.) seule :

(5.3.) Pour \inserted{tout $X\in\mathrm{Ob}\,\mathcal M$ et pour} toute partie $J$ de $\mathrm{Ss}_\pi(X)$, telle
que a) $J$ est cofinale dans $\mathrm{Ss}_\pi(X)$ et

% slanted
\margin{Au lieu de $J$ cofinale dans $\mathrm{Ss}_{\pi}(X)$, il suffit $=$ que
l'on ait $\underset{\alpha\in J}{\underrightarrow{\mathrm{lim}}}\,X_\alpha\overset{\sim}{\longrightarrow}X$,
-- la condition corr. sur $M$ devient donc plus exigeante. Mais
si les objets de $\mathrm{Filt}^{\pi}$ sont $\pi$-accessibles, cette condition
est en fait équivalente à la condition plus faible.}
b) Si \deleted{\unsure{on a}} $J'\subset J$ est une partie filtrante
de $J$, de cardinal $\leq\pi$, alors sa borne
sup. $\alpha_{J'}$ dans $\mathrm{Ss}_\pi(X)$ appartient à $J$,

\deleted{l'ens. des $J^{M\deleted{\unsure{\pi}}}$ des $\alpha$}
On a $X\in M$ \unsure{sss|ssi} $J^M\overset{\mathrm{déf}}{=}J\cap\mathrm{Ss}^{M}_\pi(X)$
est cofinal dans $J$ (ou ce qui revient
au $=$, dans $\mathrm{Ss}_\pi(X)$).

\deleted{\unsure{Compte tenu}} Si on admet déjà \unsure{établi|établie} le ``il suffit'' de
(5.1.), alors il suffit dans (5.3.) de
supposer $\deleted{\unsure{M}}X\in M$ et exiger qu'alors
$J^M$ soit cofinal dans $J$.

Mais notons le

% archive sheet number: 176
\scanpage{178}
\mspage{52}

\uline{Proposition 5.6.} Supposons $\unsure{\pi|\varpi}$ choisi de telle façon
que les objets de $\mathrm{Filt}^{\pi}$ soient $\pi$-accessibles
(p. ex. $\pi=2^c$, avec $c\geq\pi_1$\deleted{\unsure{\}}}, cf (3.2.) p. 16)).
Alors, pour une partie $M$ de $\mathrm{Ob}\,\unsure{\mathcal M|M}$, $M$ est
$\pi$-accessible \unsure{sss|ssi} elle satisfait au critère (5.3.)\textsuperscript{*}.
% slanted
\margin{Faux tel quel, vrai si l'on suppose que $M^\pi$
est stable par $\underrightarrow{\mathrm{lim}}$ filtrantes de \unsure{cart.|card.} $\leq\pi$.}

On a vu que c'est \inserted{nécessaire} (tautologiquement)\deleted{\unsure{,}}
compte tenu de prop. 5.3 et du scholie 5.5
(bas de la p. 48). Prouvons que c'est
suffisant, en établissant d'abord le

\uline{\deleted{Corollaire}}\inserted{Lemme} 5.7. Sous les conditions de
5.6. sur $\pi$, \inserted{et la condition e) sur $\mathcal M$ (p. \unsure{46|40|45}),}
pour tte représentation $\pi$-adaptée
$X=\underset{i\in I}{\underrightarrow{\mathrm{lim}}}\,X_i$ d'un objet $X$ de $\mathcal M$,
\deleted{il existe un} le sous-ensemble $I'$ \deleted{\unsure{cofinal dans}} des
$i\in I$, tels que \deleted{pour $i'\in I'$,} $X_{\unsure{i'|i}}\to X$
sont mono, est cofinal dans $I$. De plus, $I'$ est stable par \unsure{borne supérieure|bornes supérieures} \inserted{(dans $I$)} des \deleted{$\unsure{J\subset I'}$} parties filtrantes $J$ de $I'$ de cardinal $\leq\pi$.

Montrons comment ce lemme 5.7.
implique le \inserted{``il suffit'' de} prop. 5.6. \deleted{Soit} Soit donc
$X=\underset{I}{\unsure{\underline{\mathrm{lim}}|\underrightarrow{\mathrm{lim}}}}\,X_i$ une représentation $\pi$-adaptée
de $X$ ($X_i\in\mathrm{Filt}^{\pi}$, $I$ grand devant $\pi$
etc.). Il faut prouver alors, \deleted{\unsure{que}} si
5.3. est \unsure{valable}, alors $X\in M$ sss l'ens
$I_M$ est cofinal dans $I$.
\deleted{Supposons $\unsure{X\in M}$. cofinal dans $I$, appliquons
5.7, pour trouver une partie}

% archive sheet number: 177
\scanpage{179}
\mspage{53}

Supposons $X\in M$. \deleted{Appliquons (5.3.) à \unsure{l'ens.}} \inserted{Considérons le sous-ensemble}
$I'\subset I$ \deleted{\unsure{de}} décrit dans le lemme, On a une application
$I'\longrightarrow\mathrm{Ss}_{\unsure{\pi|\varpi}}(X)$, \inserted{soit $\unsure{i|i'}\longmapsto X'_i=\mathrm{Im}(X_i\hookrightarrow X)$} et on a
\[
\begin{aligned}
 X &\simeq\underset{i\in\unsure{J|I}}{\underrightarrow{\lim}}\,X_i
 \simeq\underset{i'\in I'}{\unsure{\lim|\underrightarrow{\lim}}}\,X_i
 \simeq\underset{i\in\deleted{\unsure{I}}I'}{\underrightarrow{\lim}}\,
 \underset{\substack{\Vert\\ \mathrm{Im}(X_i\hookrightarrow X)}}{X'_i},
 \deleted{\unsure{i}}\simeq\\
 &\underset{\alpha\in J}{\underrightarrow{\lim}}\,Z_\alpha ,
\end{aligned}
\]
où $J$ est l'image de $I'$ ds $\mathrm{Ss}_\pi(X)$.

\deleted{On a d'autre part} Comme les objets de $\mathrm{Filt}^{\pi}$ sont $\pi$-accessibles, on en conclut que $J$ est cofinal dans $\mathrm{Ss}_{\pi}(X)$. Par 5.3, il s'ensuit que l'ens $J_M$ \inserted{des $\alpha\in J$ tels que $Z_\alpha\in M$}
est cofinal dans $J$, \inserted{donc $\ne\varnothing$.} Donc \inserted{$I'_M$, qui est} l'image inverse \deleted{$I'_M$ de} $J_M$ dans $I'$, est en tous cas non vide.
Appliquant \deleted{\unsure{ce}} ce \unsure{$=$|--}résultat en remplaçant $I'$ par \deleted{\unsure{l'}}le sous-ensemble $I'_{\geq i'_0}$, on voit que $I'_M$ est cofinal dans $I'$, donc aussi dans $I$ puisque $I'$ est cofinal dans $I$; donc $I_M\supset I'_M$ est cofinal dans $I$, qed.

Inversement, \deleted{\unsure{prouvons}} supposons que $I_M$ soit cofinal dans $I$, et prouvons $X\in M$. La difficulté ici provient du fait que la représentation $X=\underset{i\in I_M}{\underrightarrow{\lim}}X_i$ n'est plus adaptée à $\pi$ --- \deleted{\unsure{$I_M$}} \unsure{$I_M$} n'est pas nécessairement stable \inserted{dans $I$} par bornes sup. \inserted{(dans $I$)} des parties filtrantes \inserted{de $I_M$} \unsure{$I'$} de cardinal $\leq\pi$. \unsure{En somme|Et surtout|Et souvent}, il faut

% archive sheet number: 178
